 
\section{Discussion and Conclusion}
 \vspace{-0.1in}
\noindent\textbf{Discussion.} 
Even though \ours, cannot be used for facial reenactment and hence cannot be used for harmful acts like face-swapping, it can be used for other malicious purposed like disinformation. We discourage and disapprove of any such applications which may have negative implications and strictly condone their use for positive purposes.    



\noindent\textbf{Conclusion.} In this work, we present \ours, which is able to generate high-quality lip synchronization. We pose the task to be a mouth region inpainting task and solve it by learning an audio-conditioned diffusion model. Our ablation studies show that SyncNet loss is required in our framework to introduce lip-sync while sequential adversarial loss improves both image quality and temporal consistency. Finally, extensive quantitative and qualitative results validate that our method performs better than state-of-the-art methods in terms of image quality while maintaining other metrics and also being preferred by the users. 


\noindent\textbf{Acknowledgements.} 
We would like to thank our colleagues 
Matthew Gwilliam, Archana Swaminathan, and Ahmed Taha for their feedback on this work and all the participants of the user study for their time.



